\documentclass[11pt]{amsart}
\usepackage[margin=2.6cm]{geometry}
\usepackage{amsmath,amssymb,amsthm,booktabs,enumitem}
\usepackage[hidelinks]{hyperref}
\raggedbottom
\newtheorem{theorem}{Theorem}[section]
\newtheorem{proposition}[theorem]{Proposition}
\newtheorem{lemma}[theorem]{Lemma}
\newtheorem{corollary}[theorem]{Corollary}
\newtheorem{conjecture}[theorem]{Conjecture}
\newtheorem{hypothesis}[theorem]{Hypothesis}
\theoremstyle{remark}
\newtheorem{remark}[theorem]{Remark}
\newcommand{\Z}{\mathbb Z}
\newcommand{\Q}{\mathbb Q}
\newcommand{\HH}{\mathbb H}
\newcommand{\sS}{\mathfrak S}
\newcommand{\vt}{\vartheta}
\newcommand{\dd}{\,\mathrm d}
\newcommand{\Cl}{\mathrm{Cl}}
\newcommand{\gen}{\mathrm{gen}}
\newcommand{\e}{\varepsilon}

\title[Sums of two squares and Heegner norms, VII]{Sums of two squares and Heegner norms with few prime factors, VII:\\ the Linnik problem for a single ideal class, the equidistribution\\ of the classes and the variance of the class deviation}
\author{Ruqing Chen}
\address{GUT Geoservice Inc., Montreal, Canada}
\email{ruqing@hotmail.com}
\date{\today\ (draft)}

\begin{document}
\begin{abstract}
For a prime discriminant $-D$ with class number $h>1$ and a class $C$ of binary forms $Q_C$ of discriminant $-D$ we consider $S_C(N)=\sum_{p<N}r_C(N-p)$, the number of representations $N=p+Q_C(x,y)$. The class decomposition $r_C=r_\gen+\frac wh\sum_{\psi\ne1}\bar\psi(C)c_\psi$ of Part~V splits $S_C$ into the genus part, which is the Linnik problem for $\frac wh(1*\chi_{-D})$ and has, under the hypothesis $\mathrm H(N)$ of the class-number-one case, the main term $\frac{2\pi}{\sqrt D}\sS_D(N)\pi(N)$ \emph{independently of $h$ and of the class}, and the class deviations $S_\psi(N)=\sum_{p<N}c_\psi(N-p)$, sums of the coefficients of the weight-one dihedral cusp forms $\vt_\psi$ over the shifted primes. We prove unconditionally that $\sum_{N\le X}|S_\psi(N)|^2\ll X^2\log X$, so that for all $N\le X$ outside $O(X^{1-2\eta}\log X)$ exceptions the counts of the different classes differ by at most $N^{1/2+\eta}$, and the classes are asymptotically equidistributed for all $N\le X$ but $O((\log X)^4)$; under $\mathrm H(N)$ and the dual Titchmarsh bound of Assing--Blomer--Li for $\vt_\psi$ every class has the asymptotic formula for every large $N$. Numerically ($D=23,31$, $N\le4\cdot10^6$) both classes agree with the genus formula to $10^{-4}$ and $S_\psi(N)\approx0.12\sqrt N$. The second moment is not the diagonal: $\sum_{p<N}c_\psi(N-p)^2$ is $\kappa_\psi\pi(N)$ on average, with the Rankin--Selberg residue $\kappa_\psi$ of Part~V, while the mean square of $S_\psi(N)$ is $0.41$--$0.43$ of it and rising. We explain this: by modularity the exponential sum of $\vt_\psi$ is exponentially small at the rationals $a/q$ with $q^2\ll X$ (Lemma~\ref{lem:depletion}), and these major arcs carry half of the $L^2$ mass of the prime exponential sum; we conjecture that the mean square is $\sim\frac12\kappa_\psi\pi(N)$. Under sieve weights of level $X^\theta$ on $N-p$ (towards Chen's theorem in a single class) the class deviation stays negligible for almost all $N$, unconditionally and for every $\theta<1$ (Theorem~\ref{thm:sifted}), and the depletion survives on the arcs of order $X^{1/2-\theta}$; numerically the sharp sieve by the primes below $N^{1/3}$ raises the variance ratio to $0.45$--$0.50$ and the strata $\Omega(N-p)=1,2,3$ to $0.54$--$0.81$, while the classes stay equidistributed in every stratum: there is no class-group grading of the parity obstruction, only a Chebyshev bias of lower order. The exponential sum of $c_\psi$ over the primes satisfies the Vinogradov--Siegel--Walfisz bound $X(\log X)^{-A}$ uniformly in $\alpha$ (Proposition~\ref{prop:VSW}; the twisted $L$-functions are Hecke $L$-functions of $K$ with no Siegel zero, because $h$ is odd), so the class deviation of every stratum $\Omega(N-p)=k$ and of Chen's switching sum is $O(N(\log N)^{-A})$ for almost all $N$ (Theorem~\ref{thm:strata}), and the binary Goldbach problem with one prime in a given ideal class has, for almost all even $N$, the main term $\frac2h(1-\frac{\chi(N)}{D-2})\sS_2(N)\mathrm{li}_2(N)$, the same for every class (Theorem~\ref{thm:goldbach}); for odd $N$ the representations $N=p+2q$, resp.\ $p+4q$, take its place (Corollary~\ref{cor:odd}). So for almost all $N$ and every class, $N=p+Q_C(x,y)$ with $Q_C(x,y)$ a prime, a $P_2$ or a $P_3$ according to the parity of $N$ and $\chi(2)$, unconditionally; only the problem for every individual $N$ remains.
\end{abstract}
\maketitle

\section{Setting and results}

This is Part~VII of the series. Parts~I--IV treated the theta series of the norm forms of the nine fields of class number one and the sieve for sums of two squares and Heegner norms; Part~V \cite{ChenV} passed to class number $h>1$, where the theta series of an ideal class is the genus Eisenstein series plus weight-one dihedral cusp forms, and Part~VI \cite{ChenVI} computed the spectral side of the resulting quaternary forms. Here the same class decomposition is used on the additive problem $N=p+Q_C(x,y)$: the genus part is the Linnik problem of class number one, treated in the draft \cite{ChenL}, and the new object is the class deviation, a sum of the coefficients of a weight-one cusp form over the shifted primes. Part~VIII \cite{ChenVIII} is independent of this part.


Let $D\equiv3\pmod4$ be a prime, $D>3$, $K=\Q(\sqrt{-D})$, $\chi=\chi_{-D}$, $w=2$, $h=h(-D)$ (odd), and for a class $C\in\Cl(K)$ let $r_C(n)=\#\{(x,y)\in\Z^2:\ Q_C(x,y)=n\}$ as in Part~V \cite{ChenV}. With $c_\psi(n)=\sum_{N\mathfrak a=n}\psi(\mathfrak a)$ ($\psi\in\widehat{\Cl}$, real, $c_{\bar\psi}=c_\psi$) and $r_\gen=\frac wh(1*\chi)$,
\begin{equation}\label{eq:dec}
r_C=r_\gen+\frac wh\sum_{\psi\ne1}\bar\psi(C)\,c_\psi,\qquad
S_C(N):=\sum_{p<N}r_C(N-p)=S_\gen(N)+\frac wh\sum_{\psi\ne1}\bar\psi(C)\,S_\psi(N),
\end{equation}
where $S_\gen(N)=\sum_{p<N}r_\gen(N-p)$ and $S_\psi(N)=\sum_{p<N}c_\psi(N-p)$; $\vt_\psi=\sum c_\psi(n)e(nz)\in S_1(\Gamma_0(D),\chi)$ is the dihedral cusp form of $\psi$ (Part~V, Theorem~1.1). For $h=3$, $S_{C_0}=S_\gen+\frac43S_c$, $S_{C_1}=S_\gen-\frac23S_c$ with $c=c_\psi$, and $S_{C_0}-S_{C_1}=2S_c$: for $D=23$, $S_c(N)$ is half the difference between the numbers of representations $N=p+x^2+xy+6y^2$ and $N=p+2x^2+xy+3y^2$.

\subsection{The genus part}
Since $r_\gen=\frac wh(1*\chi)$ and the draft \cite{ChenL} treats $r_D=w_D(1*\chi)$ for $h=1$, its proof applies verbatim with $w_D$ replaced by $\frac wh$, and the constant becomes $\frac whL(1,\chi)=\frac{2\pi}{\sqrt D}$ by the class number formula:

\begin{proposition}\label{prop:genus}
Under the hypothesis $\mathrm H(N)$ of \cite[Hypothesis~1.1]{ChenL} (the level $\sqrt N(\log N)^A$ of Bombieri--Friedlander--Iwaniec for the single residue class $N$),
\[
S_\gen(N)=\frac{2\pi}{\sqrt D}\,\sS_D(N)\,\pi(N)+O\Bigl(N\Bigl(\frac{\log\log N}{\log N}\Bigr)^2\Bigr),\]\[\sS_D(N)=\Bigl(1-\frac{\chi(N)}{D-1}\Bigr)\prod_{\ell\mid N}\Bigl(1-\frac{\chi(\ell)}\ell\Bigr)\prod_{\ell\nmid N}\Bigl(1+\frac{\chi(\ell)}{\ell(\ell-1)}\Bigr).
\]
\end{proposition}
The main term is the same as for the principal form of a field of class number one with the same $D$: the class number cancels between the genus average and the class number formula, exactly as the leading coefficient of the scalloped norms in Part~V. Every class has the local densities of the genus (there is one genus), so the singular series is the same for all classes, and \eqref{eq:dec} reduces the Linnik problem for a single class to the size of $S_\psi(N)$.

\subsection{The class deviation}
\begin{lemma}[Wilton-type bound]\label{lem:wilton}
Let $f=\sum a(n)e(nz)\in S_1(\Gamma_0(D),\chi)$. Then, uniformly in $\alpha\in\mathbb R$ and $Y\ge2$, $\sum_{n\le Y}a(n)e(n\alpha)\ll_f\sqrt Y\log Y$, and $\sum_na(n)e^{-n/Y}e(n\alpha)=f(\alpha+\frac i{2\pi Y})\ll_f\sqrt Y$.
\end{lemma}
\begin{proof}
The function $y^{1/2}|f(z)|$ is $\Gamma_0(D)$-invariant ($|\chi|=1$) and tends to $0$ at both cusps, so it is bounded on $\HH$: $|f(x+iy)|\le C_fy^{-1/2}$. This is the second statement. For the first, with $y=1/Y$, $\sum_{n\le Y}a(n)e(n\alpha)=\int_0^1f(\alpha+\beta+iy)K(\beta)\dd\beta$, $K(\beta)=\sum_{m\le Y}e^{2\pi m/Y}e(-m\beta)$; by partial summation against the Dirichlet kernels, whose $L^1$ norms are $\ll\log Y$, and since $e^{2\pi m/Y}$ is increasing and $\le e^{2\pi}$, $\|K\|_1\ll\log Y$.
\end{proof}

\begin{theorem}\label{thm:almostall}
For every $\psi\ne1$ and $X\ge2$,
\[
\sum_{N\le X}|S_\psi(N)|^2\ll_DX^2\log X .
\]
Consequently, for every $\eta>0$, all $N\le X$ outside a set of $O_D(X^{1-2\eta}\log X)$ elements satisfy $|S_C(N)-S_\gen(N)|\le N^{1/2+\eta}$ for every class $C$; and for every $\e>0$, all $N\le X$ outside a set of $O_{D,\e}((\log X)^4)$ elements satisfy $|S_C(N)-S_\gen(N)|\le\e N/\log N$ for every class.
\end{theorem}
\begin{proof}
For $N\le X$, $S_\psi(N)$ is the coefficient of $e(N\alpha)$ in $P(\alpha)F(\alpha)$ with $P(\alpha)=\sum_{p\le X}e(p\alpha)$ and $F(\alpha)=\sum_{1\le m\le X}c_\psi(m)e(m\alpha)$, since $p<N$ and $m=N-p$ are both $\le X$. By Parseval and Lemma~\ref{lem:wilton},
\[
\sum_{N\le X}|S_\psi(N)|^2\le\int_0^1|P(\alpha)F(\alpha)|^2\dd\alpha\le\|F\|_\infty^2\int_0^1|P|^2=\|F\|^2_\infty\,\pi(X)\ll X(\log X)^2\frac X{\log X}.
\]
By \eqref{eq:dec} the same bound holds for $S_C-S_\gen$. The number of $N\in(X/2,X]$ with $|S_C(N)-S_\gen(N)|>(X/2)^{1/2+\eta}$ is $\ll X^2\log X/X^{1+2\eta}$, and with $>\e X/(2\log X)$ it is $\ll(\log X)^3/\e^2$; summing over dyadic ranges gives the two counts.
\end{proof}

Thus the equidistribution of the representations among the classes holds unconditionally for almost all $N$, with a power saving, while the asymptotic formula itself, which is a statement about the genus part, needs $\mathrm H(N)$ already for $h=1$. For individual $N$ the class deviation needs an input of a different kind, a bound for the coefficients of a cusp form along $N-p$; this is the dual problem of Titchmarsh type.

\begin{hypothesis}[$\mathrm{T}_\psi$]\label{hyp:T}
$S_\psi(N)=o(N/\log N)$ as $N\to\infty$.
\end{hypothesis}

Assing, Blomer and Li \cite[Theorem~1.3]{ABL} prove, for a holomorphic or Maa\ss{} cuspidal newform for $\Gamma_0(M)$ with Hecke eigenvalues $\lambda(n)$, $\sum_{Y<n\le X}\Lambda(n)\lambda(f-n)\ll X^{1-\delta}$ for $f>X$, $X+f\asymp X$ (the case $\sigma=-1$ of their theorem, which they describe as the cuspidal analogue of the dual sum $\sum_{p<N}\tau(N-p)$ and which was out of reach before their work). With $f=N$, $Y=0$, $X=N-1$ their theorem bounds $\sum_{n<N}\Lambda(n)\lambda(N-n)$ by $N^{1-\delta}$, and partial summation gives $\mathrm T_\psi$ with a power saving, provided (i) the theorem extends to the weight-one forms $\vt_\psi$, which have the nebentypus $\chi_{-D}$ (their statement is for $\Gamma_0(M)$ with trivial character), and (ii) it holds for the truncations $n\le t$, $N/(\log N)^{A}\le t\le N$, i.e.\ with $X+f\ll X(\log X)^A$ instead of $X+f\asymp X$. Remark~\ref{rem:ABL} records what (i) requires; (ii) we have not checked. We therefore state:

\begin{theorem}\label{thm:every}
Under $\mathrm H(N)$ and $\mathrm T_\psi$ for all $\psi\ne1$, for every class $C$,
\[
S_C(N)=\frac{2\pi}{\sqrt D}\,\sS_D(N)\,\pi(N)\,(1+o(1)),
\]
the same main term for every class; in particular every large $N$ is $p+Q_C(x,y)$ for every class $C$, and the representations are equidistributed among the classes.
\end{theorem}
\begin{proof}
Proposition~\ref{prop:genus}, \eqref{eq:dec}, $\mathrm T_\psi$, and $\sS_D(N)\gg1$ \cite[Theorem~1.2]{ChenL}.
\end{proof}

\section{Numerics}
Table~\ref{tab:ratios} compares $S_{C_0}$, $S_{C_1}$ with the main term of Proposition~\ref{prop:genus} for $D=23$ and $31$, on $400$ random $N$ in each range (script \texttt{linnik\_class.py}, exact counts by FFT from the representation numbers). The class deviation is of square-root size, $\mathrm{rms}\,S_c(N)\approx0.12\sqrt N$ (a fit of $\log\mathrm{rms}$ against $\log N$ gives the exponent $0.46$ for both $D$), and it has no bias: its mean over $N\in[2\cdot10^6,4\cdot10^6)$ is $0.000\sqrt N$ separately for $\chi(N)=1,-1,0$.

\begin{table}[htbp]
\centering\footnotesize\setlength{\tabcolsep}{5pt}
\begin{tabular}{r|rrrr|rrrr}
\toprule
 & \multicolumn{4}{c|}{$D=23$} & \multicolumn{4}{c}{$D=31$}\\
$N$ & $S_{C_0}/\mathrm{main}$ & $S_{C_1}/\mathrm{main}$ & $\frac{\mathrm{rms}\,S_c}{\sqrt N}$ & $\frac{E\,S_c^2}{\kappa_D\pi(N)}$ & $S_{C_0}/\mathrm{main}$ & $S_{C_1}/\mathrm{main}$ & $\frac{\mathrm{rms}\,S_c}{\sqrt N}$ & $\frac{E\,S_c^2}{\kappa_D\pi(N)}$\\
\midrule
$[10^5,2\cdot10^5)$ & $1.0001$ & $1.0001$ & $0.126$ & $0.405$ & $0.9997$ & $0.9998$ & $0.131$ & $0.414$\\
$[3\cdot10^5,6\cdot10^5)$ & $1.0002$ & $0.9999$ & $0.125$ & $0.414$ & $1.0000$ & $1.0003$ & $0.126$ & $0.422$\\
$[10^6,2\cdot10^6)$ & $1.0005$ & $1.0003$ & $0.119$ & $0.423$ & $1.0001$ & $1.0000$ & $0.119$ & $0.430$\\
$[2\cdot10^6,4\cdot10^6)$ & $0.9999$ & $1.0000$ & $0.113$ & $0.427$ & $1.0002$ & $0.9999$ & $0.119$ & $0.432$\\
\bottomrule
\end{tabular}
\medskip
\caption{The Linnik problem for the two classes of $D=23$ ($x^2+xy+6y^2$, $2x^2+xy+3y^2$) and $D=31$ ($x^2+xy+8y^2$, $2x^2+xy+4y^2$): the counts against $\frac{2\pi}{\sqrt D}\sS_D(N)\pi(N)$ (means over $400$ random $N$; standard deviations $0.002$--$0.03$), the class deviation $S_c=\frac12(S_{C_0}-S_{C_1})$, and its mean square over the whole range against $\kappa_D\pi(N)$, $\kappa_{23}=0.421799$, $\kappa_{31}=0.430026$ (\texttt{linnik\_class\_variance.py}). The diagonal $\sum_{p<N}c(N-p)^2$ is $\kappa_D\pi(N)$ to $10^{-4}$ in every range; for $D=23$ and $N\in[8\cdot10^6,1.6\cdot10^7)$ the last column is $0.419$ for $\chi(N)=1$, $0.448$ for $\chi(N)=-1$, $0.433$ for $\chi(N)=0$.}
\label{tab:ratios}
\end{table}

\section{The variance of the class deviation}
The diagonal of $S_\psi(N)^2$ is $\sum_{p<N}c_\psi(N-p)^2$, and since $\sum_{m\le x}c_\psi(m)^2=\kappa_\psi x+O(x^{4/5})$ with the residue $\kappa_\psi$ of Part~V, Theorem~1.3(iii), its average over $N\le X$ is $\kappa_\psi\pi(N)$ (Table~\ref{tab:ratios}: to $10^{-4}$). The mean square of $S_\psi(N)$ is much smaller, $0.41$--$0.43\,\kappa_\psi\pi(N)$, increasing slowly with $N$. The off-diagonal correlations $\sum_{p\ne p'}c_\psi(N-p)c_\psi(N-p')$ are therefore negative on average and of the size of the diagonal. The mechanism is visible on the circle.

\begin{lemma}[Depletion at the major arcs]\label{lem:depletion}
Let $f\in S_1(\Gamma_0(D),\chi)$, $F_X(\alpha)=\sum_na(n)e^{-n/X}e(n\alpha)=f(\alpha+\frac i{2\pi X})$, and $(a,q)=1$. Then
\[
|F_X(a/q)|\ll_f\frac Xq\exp\Bigl(-\frac{4\pi^2X}{Dq^2}\Bigr)\quad(D\nmid q),\qquad |F_X(a/q)|\ll_f\frac Xq\exp\Bigl(-\frac{4\pi^2X}{q^2}\Bigr)\quad(D\mid q).
\]
\end{lemma}
\begin{proof}
Let $\gamma=\left(\begin{smallmatrix}a&b\\q&d\end{smallmatrix}\right)\in\mathrm{SL}_2(\Z)$ and $z=a/q+iy=\gamma w$, so that $w=-d/q+i/(q^2y)$ and $|qw+d|=1/(qy)$. Then $f(z)=(qw+d)\,(f|_1\gamma)(w)$. If $D\mid q$, $\gamma\in\Gamma_0(D)$ and $f|_1\gamma=\chi(d)f$; if $D\nmid q$, $\gamma$ maps $\infty$ to a cusp equivalent to $0$, and $f|_1\gamma$ is the expansion of $f$ at $0$, $\sum_{n\ge1}b(n)e(nw/D)$ ($f|W_D=-if$ for the dihedral forms, Part~V, Theorem~1.1). In both cases $|f|_1\gamma(w)|\ll e^{-2\pi\operatorname{Im}w/D'}$ with $D'=D$ or $1$, and $\operatorname{Im}w=1/(q^2y)=2\pi X/q^2$.
\end{proof}

So $F_X$ vanishes to exponential precision at every $a/q$ with $q^2\le\delta X/\log X$, whereas its mean square over $[0,1]$ is $\sum c_\psi(n)^2e^{-2n/X}\sim\frac12\kappa_\psi X$. Numerically, with $\rho(q)$ the mean of $|F_X(a/q)|^2$ over $(a,q)=1$ divided by this mean square ($D=23$, $X=3\cdot10^5$, \texttt{linnik\_class\_arcs.py}): $\rho(q)<10^{-4}$ for $q^2<0.26X$, then $0.008$, $0.04$, $0.12$, $0.50$ at $q^2/X=0.45$, $0.59$, $0.77$, $1.34$, an overshoot to $1.50$ at $q^2/X\approx7$, and $\rho\to1$ (e.g.\ $1.07$ at $q^2/X=36$). The whole profile, overshoot included, is given by the coefficients $c_\psi(n)^2$ with $n$ of size $q^2/X$:

\begin{proposition}[Scaling profile of the depletion]\label{prop:profile}
Let $\Phi_D(Y)=\frac1{\kappa_\psi Y}\sum_{n\ge1}c_\psi(n)^2e^{-n/Y}$, and let $\rho(q)=\frac1{\varphi(q)}\sum_{(a,q)=1}|F_X(a/q)|^2\big/\kappa_\psi\sum_me^{-2m/X}$ for $F_X$ as in Lemma~\ref{lem:depletion} with $f=\vt_\psi$. Then
\[
\rho(q)=\Phi_D(Y)\bigl(1+O(X^{-1})\bigr)+O\Bigl(\frac{(1+Y)^2\log^2(2+Y)}{\varphi(q)}\Bigr),\qquad Y=\begin{cases}\dfrac{Dq^2}{8\pi^2X},&D\nmid q,\\[2mm]\dfrac{q^2}{8\pi^2X},&D\mid q.\end{cases}
\]
\end{proposition}
\begin{proof}
Let $D\nmid q$, $\gamma=\left(\begin{smallmatrix}a&b\\q&d\end{smallmatrix}\right)\in\mathrm{SL}_2(\Z)$, and choose $u,v$ with $uq-Dva=1$, so that $\delta=\left(\begin{smallmatrix}u&a\\Dv&q\end{smallmatrix}\right)\in\Gamma_0(D)$ and
\[
\delta^{-1}\gamma=\begin{pmatrix}0&-1\\1&k\end{pmatrix}=W_D\cdot\frac1D\begin{pmatrix}1&k\\0&D\end{pmatrix},\qquad k=ud-Dvb.
\] With $\vt_\psi|_1W_D=-i\vt_\psi$ this gives $|\vt_\psi|_1\gamma(w)|=D^{-1/2}|\vt_\psi(\frac{w+k}D)|$. At $z=\frac aq+i\eta$, $\eta=\frac1{2\pi X}$, one has $w=\gamma^{-1}z=-\frac dq+\frac i{q^2\eta}$ and $kq-d=Dv$, so
\[
|F_X(a/q)|=\frac1{q\eta\sqrt D}\Bigl|\sum_nc_\psi(n)\,e\Bigl(-\frac{n\,\overline{aD}}q\Bigr)e^{-2\pi n/(Dq^2\eta)}\Bigr|,
\]
since $v\equiv-\overline{Da}\pmod q$. Averaging over $(a,q)=1$ gives, with the Ramanujan sum $c_q$,
\[
\frac1{Dq^2\eta^2}\sum_{n_1,n_2}c_\psi(n_1)c_\psi(n_2)e^{-2\pi(n_1+n_2)/(Dq^2\eta)}\frac{c_q(n_1-n_2)}{\varphi(q)}.
\] The diagonal is $\frac{X}{2Y}\sum_nc_\psi(n)^2e^{-n/Y}$, and $\kappa_\psi\sum_me^{-2m/X}=\frac{\kappa_\psi X}2(1+O(X^{-1}))$; off the diagonal $|c_q(h)|\le(h,q)\le|h|$ and $|c_\psi(n)|\le\tau(n)$ give the error term. For $D\mid q$, $\gamma\in\Gamma_0(D)$, $|\vt_\psi|_1\gamma|=|\vt_\psi|$, and the same computation runs with the width $1$ in place of $D$.
\end{proof}

At $X=3\cdot10^5$, $D=23$ (\texttt{linnik\_class\_arcs.py}) the two sides agree to $\pm0.002$ at all $q$, e.g.\ $\rho=0.0084$, $0.1245$, $0.5039$, $1.5026$, $1.0720$ against $\Phi_D=0.0084$, $0.1245$, $0.5047$, $1.5022$, $1.0708$ at $q^2/X=0.45$, $0.77$, $1.34$, $6.93$, $36.0$, and the arcs with $D\mid q$, which by the proposition stay depleted up to $q^2/X\approx8\pi^2\cdot0.1=7.9$ instead of $0.34$, follow $\Phi_D(q^2/8\pi^2X)$ just as closely ($\rho=0.0011$, $0.0577$, $0.5186$ at $q^2/X=7.9$, $14.6$, $31.2$). The overshoot is the excess of the first coefficients over their mean, $c_\psi(1)^2=c_\psi(2)^2=c_\psi(3)^2=1$ against $\kappa_{23}=0.42$: $\max\Phi_{23}=1.503$ at $Y=1.90$, i.e.\ $q^2/X=6.5$. The logarithmic integral of the profile is a Rankin--Selberg constant: if $\sum_nc_\psi(n)^2n^{-s}=\frac{\kappa_\psi}{s-1}+c_0+O(s-1)$, then, from $\int_0^\infty\sum_nc_\psi(n)^2e^{-nt}t^{s-1}\dd t=\Gamma(s)\sum_nc_\psi(n)^2n^{-s}$ at $s=1$,
\[
\int_0^\infty\bigl(\mathbf 1_{Y>1}-\Phi_D(Y)\bigr)\frac{\dd Y}Y=\gamma-\frac{c_0}{\kappa_\psi};
\]
for $D=23$, $c_0=1.03035$ ($=\lim(\sum_{n\le x}c_\psi(n)^2/n-\kappa_\psi\log x)$) and both sides are $-1.866$ ($-1.868$ by quadrature). With $\sum_{q\le Q,\,D\nmid q}\frac{\mu^2(q)}{\varphi(q)}\sim\frac{D-1}D(\log Q+1.3325+\frac{\log D}D)$ and the analogue for $D\mid q$, this evaluates the deficit of the model below in closed form, $\sum_q\frac{\mu^2(q)}{\varphi(q)}(1-\rho(q))=\frac12\log X+\frac12\log8\pi^2-\frac{D-1}{2D}\log D+1.3325+\frac12(\gamma-\frac{c_0}{\kappa_\psi})+o(1)$, i.e.\ $7.39$ at $X=3\cdot10^5$, $D=23$.

The prime exponential sum does the opposite: $|P(a/q+\beta)|^2\approx\frac{\mu(q)^2}{\varphi(q)^2}|T(\beta)|^2$ with $\int|T|^2\approx X/(\log X)^2$, so the arcs at the $\varphi(q)$ fractions $a/q$ carry the $L^2$ mass $\frac{\mu(q)^2}{\varphi(q)}\frac X{(\log X)^2}$, and the arcs with $q\le Q$ carry $\frac{X}{(\log X)^2}\sum_{q\le Q}\frac{\mu(q)^2}{\varphi(q)}=\frac X{(\log X)^2}(\log Q+1.3325+o(1))$, out of the total $\pi(X)\approx X/\log X$. (The mean square of the primes in the residue classes modulo $q$, summed over $q\le Q$, is asymptotically known for $Q\le X$ by the theorem of Barban--Davenport--Halberstam in the form of Montgomery \cite{Montgomery} and Hooley \cite{Hooley}; it is the quantity that measures how much of the mass of $|P|^2$ the major arcs carry.) The arcs with $q\le\sqrt X$ carry half the mass, and on them $F$ is exponentially small. Writing $\int|P|^2|F|^2$ as the diagonal minus the deficit on the major arcs gives the prediction
\[
\frac{E\,|S_\psi(N)|^2}{\kappa_\psi\pi(N)}\approx1-\frac{\sum_q\mu(q)^2\varphi(q)^{-1}(1-\rho(q))}{\log X+1.3325}=\frac{\frac12\log X+O(1)}{\log X+1.3325},
\]
which for $X=3\cdot10^5$ evaluates to $1-7.449/13.94=0.466$ ($0.470$ with the closed form $7.39$ after Proposition~\ref{prop:profile}), against the measured $0.41$--$0.42$ in the ranges around $3\cdot10^5$. The model is only qualitative at these sizes, and the reason is the normalisation, not the profile, which is exact (Proposition~\ref{prop:profile}). The arc masses $\frac{\mu^2(q)}{\varphi(q)}\int|T|^2$ have to be compared with $\int|P|^2$, and $\int|P|^2/\int|T|^2$ is $\log X-\log2-\gamma+o(1)$ for the smoothed sums ($8.23$, $9.64$, $10.86$ exactly at $X=3\cdot10^4$, $10^5$, $3\cdot10^5$), not $\log X+1.3325=\sum_{q\le X}\mu^2(q)/\varphi(q)$. With the correct normalisation the model gives $0.24$, $0.29$, $0.32$ at $\beta=0$ and $0.29$, $0.33$ with the $\beta$-integration of \texttt{linnik\_class\_smooth.py}, against the exact $0.393$, $0.403$. The model now has too much depletion, as it had too little before. The arcs overlap in the transition zone $q\asymp\sqrt X$, where the spacing of the fractions is of the order of the width $1/X$, and summed over $q\le X$ the arc masses exceed $\int|P|^2$ by the factor $(\log X+1.33)/(\log X-1.27)$. The arc model therefore does not determine the $O(1/\log X)$ term, whichever normalisation is used. The leading $\frac12$ and the reduction of Section~\ref{subsec:reduction}, which uses only disjoint arcs, are unaffected. With the old normalisation, Proposition~\ref{prop:profile} gives the model constant $b_{\mathrm{model}}=-\frac12\log8\pi^2+\frac{D-1}{2D}\log D-\frac12(\gamma-\frac{c_0}{\kappa_\psi})=0.25$ for $D=23$; the measured $b$ below is $-0.52$. For the sharp sums of Table~\ref{tab:ratios} the correct normalisation is $\int|P|^2/\int|T|^2=\pi(N)/\sum_{2\le n\le N}(\log n)^{-2}=\log N-1+O(1/\log N)$, and replacing $\log N+1.3325$ by $\log N-1$ shifts the model constant by $-\frac12(1+1.3325)=-1.17$, to $b_{\mathrm{model}}=-0.92$. Asymptotically the model therefore errs by $0.4$ on the other side. At these sizes the $O(1/\log N)$ terms are still $0.3$--$0.5$ (e.g.\ $11.29$ against $\log N-1=11.61$ at $N=3\cdot10^5$). For the smoothed variance $V(X)=\sum_Ne^{-2N/X}S_\psi(N)^2$, which equals $\int_0^1|P|^2|F|^2$ exactly for $P=\sum_pe^{-p/X}e(p\alpha)$ and $F=F_X$, the exact ratio to the diagonal is $0.393$ and $0.403$ at $X=3\cdot10^4$ and $10^5$ (\texttt{linnik\_class\_smooth.py}), while the same model with $|F(a/q+\beta)|^2$ integrated against the width $|T(\beta)|^2$ of the arcs gives $0.499$ and $0.498$: the integration over $\beta$ weakens the depletion ($\rho=0.21$ instead of $<10^{-4}$ at $q^2/X=0.1$, because $\operatorname{Im}w$ in Lemma~\ref{lem:depletion} drops by the factor $1+4\pi^2X^2\beta^2$ across the arc), and the model then overestimates $V$ by $0.1$, an $O(1/\log X)$ error it does not capture (the normalisation of the arc masses, see above: correctly normalised the same model underestimates $V$ by $0.07$--$0.1$; and the local structure at $D$: the measured ratio is $0.412$ for $\chi(N)=1$ and $0.442$ for $\chi(N)=-1$ around $3\cdot10^5$, the diagonal itself being $0.955$ and $1.045$ times $\kappa_D\pi(N)$ because $m=N-p$ avoids the class $N$ modulo $D$). The leading constant $\frac12$ does not depend on these details: the transition region and the overshoot change only the $O(1)$, and fitting $V(\log N+1.3325)=a\log N+b$ to the four ranges of Table~\ref{tab:ratios} gives $a=0.52$ for $D=23$, and the constant $b=V(\log N+1.3325)-\frac12\log N$ is $-0.59$, $-0.57$, $-0.53$, $-0.52$, $-0.52$ at $N\approx1.5\cdot10^5$, $4.5\cdot10^5$, $1.5\cdot10^6$, $3\cdot10^6$, $1.2\cdot10^7$: stable, as the limit $\frac12$ requires (a different limit $a$ would make $b$ drift linearly in $\log N$); the two smoothed values give the slope $a=0.50$ as well.

\begin{conjecture}\label{conj:var}
For every $\psi\ne1$, as $X\to\infty$,
\[
\sum_{X<N\le2X}|S_\psi(N)|^2\sim\frac12\,\kappa_\psi\sum_{X<N\le2X}\pi(N),
\]
and the same with the exponent $\frac12$ for the coefficients of any fixed holomorphic cusp form (with $\kappa_\psi$ replaced by its Rankin--Selberg residue, normalised to weight one): the class deviation of the Linnik problem has half the variance of a sum of independent terms.
\end{conjecture}

\begin{remark}
(a) The half is a property of the circle, not of the class group: it comes from the modularity of $f$, which makes $F$ small at the rationals of height up to $\sqrt X$, and from the logarithmic distribution of the mass of $|P|^2$ over the denominators. Section~\ref{subsec:reduction} reduces the conjecture to two statements about the primes alone and proves everything else. (b) For $h=1$ the analogue of $S_\psi$ does not exist (there is no cusp form), and for composite $D$ with several genera the genus characters give Eisenstein series, for which $F$ is large at the major arcs and the class deviation has a main term (Part~V, Section~7): the class equidistribution of Theorem~\ref{thm:almostall} is a phenomenon of the single genus. (c) For $h=3$, Theorem~\ref{thm:almostall} and Conjecture~\ref{conj:var} say that $\#\{N=p+x^2+xy+6y^2\}-\#\{N=p+2x^2+xy+3y^2\}=2S_c(N)$ is typically of size $\sqrt{2\kappa_{23}N/\log N}=0.92\sqrt{N/\log N}$, against a main term of size $1.31\,N/\log N$ for each.
\end{remark}

\subsection{What is proved about the variance}\label{subsec:reduction}
We work with the smoothed objects $P(\alpha)=\sum_pe^{-p/X}e(p\alpha)$, $F(\alpha)=F_X(\alpha)=\sum_mc_\psi(m)e^{-m/X}e(m\alpha)$ and
\begin{gather*}
V(X):=\sum_Ne^{-2N/X}S_\psi(N)^2=\int_0^1|P|^2|F|^2,\\
\mathrm{Diag}(X):=\sum_Ne^{-2N/X}\sum_{p<N}c_\psi(N-p)^2=\int_0^1|P|^2\int_0^1|F|^2,
\end{gather*}
both identities exact ($e^{-p/X}e^{-m/X}=e^{-N/X}$; the second is Parseval for each factor), so $V/\mathrm{Diag}$ is the correlation coefficient of $|P|^2$ and $|F|^2$ on the circle, and $\mathrm{Diag}\sim\frac{\kappa_\psi}2X\cdot\frac X{2\log X}$. For $B\ge3$ let
\begin{gather*}
Q=X^{1/2}(\log X)^{-B},\qquad K=(\log X)^{(B-1)/2},\\
\mathfrak M=\bigcup_{q\le Q}\ \bigcup_{(a,q)=1}\Bigl[\frac aq-\frac KX,\frac aq+\frac KX\Bigr],\qquad \mathfrak m=[0,1)\setminus\mathfrak M
\end{gather*}
(the arcs are disjoint, $2K/X<Q^{-2}$).

\begin{theorem}\label{thm:reduction}
$\sup_{\alpha\in\mathfrak M}|F(\alpha)|\ll_{A}X^{-A}$ for every $A$; hence $V(X)=\int_{\mathfrak m}|P|^2|F|^2+O_A(X^{-A})$.
\end{theorem}
\begin{proof}
Let $\alpha=\frac aq+\beta\in\mathfrak M$, $z=\alpha+iy$, $y=\frac1{2\pi X}$, so that $F(\alpha)=f(z)$, and $\gamma=\left(\begin{smallmatrix}a&b\\q&d\end{smallmatrix}\right)\in\mathrm{SL}_2(\Z)$, $z=\gamma w$. As in Lemma~\ref{lem:depletion}, $|f(z)|=|f|_1\gamma(w)|/(q|\beta+iy|)\le|f|_1\gamma(w)|/(qy)$ with
\[
\operatorname{Im}w=\frac{y}{q^2(\beta^2+y^2)}\ge\frac{2\pi X}{q^2(1+4\pi^2K^2)}\ge\frac{2\pi(\log X)^{2B}}{1+4\pi^2(\log X)^{B-1}}\gg(\log X)^{B+1},
\]
and $|f|_1\gamma(w)|\ll e^{-2\pi\operatorname{Im}w/D}$ for $\operatorname{Im}w\ge1$ (the expansion at the cusp $\gamma\infty$ has width $1$ or $D$ and polynomially bounded coefficients). So $|F(\alpha)|\ll X\exp(-c(\log X)^{B+1})$, and $\int_{\mathfrak M}|P|^2|F|^2\le\sup_{\mathfrak M}|F|^2\int_0^1|P|^2$.
\end{proof}

Thus the major arcs of order $X^{1/2-\e}$ contribute nothing to the variance, whatever the primes do there. Conjecture~\ref{conj:var} (in the smoothed form $V\sim\frac12\mathrm{Diag}$) follows from the two statements
\begin{align*}
&\textup{(M)}\quad \int_{\mathfrak M}|P|^2=\Bigl(\frac12+o(1)\Bigr)\int_0^1|P|^2\qquad(\text{for some fixed }B,\ \text{then }o(1)=O(\log\log X/\log X)),\\
&\textup{(D)}\quad \int_{\mathfrak m}|P|^2|F|^2=(1+o(1))\int_{\mathfrak m}|P|^2\cdot\int_0^1|F|^2,
\end{align*}
and conversely (M) and the conjecture imply (D). Statement (M) is the Hardy--Littlewood prediction for the mass of the arcs: with $P^*(\frac aq+\beta)=\frac{\mu(q)}{\varphi(q)}T(\beta)$, $T(\beta)=\sum_{n\ge2}\frac{e^{-n/X}}{\log n}e(n\beta)$, one has $\int_{\mathfrak M}|P^*|^2=\sum_{q\le Q}\frac{\mu(q)^2}{\varphi(q)}\int_{|\beta|\le K/X}|T|^2=(\log Q+O(1))\frac X{2(\log X)^2}(1+O(K^{-1}))$, which is $(\frac12-\frac{B\log\log X}{\log X}+O(\frac1{\log X}))\int_0^1|P|^2$. A proof of (M) needs $\int_{\mathfrak M}|P-P^*|^2=o(X/\log X)$ for arcs of order $X^{1/2-\e}$. By Gallagher's lemma and the expansion of $e(an/q)$ in characters this is a mean value $\sum_{r\le Q}\frac r{\varphi(r)}\log\frac Qr\sum_{\chi\bmod r}^*|\psi(t+H,\chi)-\psi(t,\chi)|^2\ll H^2(\log X)^{-A}$ over intervals of length $H=X/2K$: a Barban--Davenport--Halberstam estimate in short intervals with the weight $r$ instead of $1/\varphi(r)$. The large sieve gives $(H+Q^2)H\log X$, the trivial size, and the Bombieri--Vinogradov theorem, Hooley's form of the Barban--Davenport--Halberstam theorem and the mean-square Hardy--Littlewood theorem of Lavrik and Mikawa for prime pairs all lose a power of $Q$ against this weight; under GRH the pointwise bound $P-P^*\ll(qX)^{1/2}(\log X)^2(1+X|\beta|)$ gives (M) for arcs of order $X^{1/3-\e}$ only, i.e.\ the mass $\frac13$. So (M) is open, unconditionally and under GRH. Statement (D) is a genuine correlation statement between the primes and the cusp form on the minor arcs, of the kind that the pair correlation hypothesis supplies in the theorem of Goldston and Montgomery \cite{GM} on the variance of the primes in short intervals, $\int_1^X(\psi(x+\delta x)-\psi(x)-\delta x)^2\dd x\sim\frac12\delta X^2\log\frac1\delta$: there the deficit $\log\frac1\delta$ against $\log X$ is the mass of the arcs of order $\delta^{-1}$ on which the short interval kernel vanishes, as here it is the mass of the arcs of order $\sqrt X$ on which $F$ vanishes, and the analogue of $\delta=X^{-1/2}$ gives the $\frac12$.

Numerically (\texttt{linnik\_class\_split.py}, $D=23$, the arcs of order $c\sqrt X$ with $K=1$; the $c$ below are fixed, so the arcs include part of the transition zone of $\rho$):
\begin{center}\footnotesize
\begin{tabular}{r|c|cc|ccc}
\toprule
$X$ & $V/\mathrm{Diag}$ & $\int_{\mathfrak M}|P|^2/\int|P|^2$, $c=\frac14$ & HL value & (D)-ratio, $c=\frac18$ & $c=\frac14$ & $c=\frac12$\\
\midrule
$5\cdot10^4$ & $0.397$ & $0.4870$ & $0.4899$ & -- & $0.771$ & $0.866$\\
$2\cdot10^5$ & $0.408$ & $0.4885$ & $0.4856$ & $0.709$ & $0.794$ & $0.879$\\
$6\cdot10^5$ & $0.415$ & $0.4855$ & $0.4797$ & $0.727$ & $0.803$ & $0.883$\\
\bottomrule
\end{tabular}
\end{center}
The share of $V$ on $\mathfrak M$ is $0.0000$ for $c=\frac18$ and $0.003$--$0.005$ for $c=\frac14$ (Theorem~\ref{thm:reduction}); the mass of $\mathfrak M$ agrees with the Hardy--Littlewood value to $0.006$ (statement (M) at these sizes); and the ratio in (D), $\int_{\mathfrak m}|P|^2|F|^2/(\int_{\mathfrak m}|P|^2\int|F|^2)$, is $0.71$--$0.88$ and increases with $X$ for each fixed $c$, as it must if the deficit is carried by the transition zone $c\sqrt X\le q\le C\sqrt X$, whose share of the mass of $|P|^2$ is $O(1/\log X)$. Fitting $V/\mathrm{Diag}\cdot(\log X+1.3325)=a\log X+b$ to the three values gives $a=0.504$ and $0.505$ on the two intervals.

\subsection{Several modes: class number five}\label{subsec:h5}
For $h=5$ there are two pairs $\{\psi_j,\overline\psi_j\}$, $\psi_j(C_1)=e(j/5)$, and two real sequences $c_j=c_{\psi_j}$ with $S_C-S_\mathrm{gen}=\frac25\sum_{j=1,2}a_j(C)S_j$, $a_j(C_1^k)=2\cos\frac{2\pi jk}5$, i.e.\ $a=(2,2)$, $(\phi^{-1},-\phi)$, $(-\phi,\phi^{-1})$ for the three forms ($\phi$ the golden ratio). The cross diagonal $\sum_{p<N}c_1(N-p)c_2(N-p)$ is $o(\pi(N))$: $\sum_nc_1(n)c_2(n)n^{-s}$ is, up to a factor holomorphic at $s=1$, $L(s,\psi_1\psi_2)L(s,\psi_1\overline\psi_2)$ with both characters non-trivial, so it has no pole. The smoothed covariance is $\int|P|^2\operatorname{Re}F_1\overline{F_2}$, Theorem~\ref{thm:reduction} applies to each $F_j$, and the natural extension of Conjecture~\ref{conj:var} is that $E\,S_1S_2=o(\pi(N))$ (on the minor arcs $F_1\overline{F_2}$ has mean $o(1)$ times $\int|F_1||F_2|$ and should decorrelate from $|P|^2$ as in (D)), so that
\[
E\,(S_C(N)-S_\mathrm{gen}(N))^2\sim\frac2{25}\bigl(a_1(C)^2\kappa_1+a_2(C)^2\kappa_2\bigr)\pi(N):
\]
the variances of the classes are fixed combinations of the two residues with golden-ratio weights, e.g.\ in the ratio $\frac{\phi^{-2}\kappa_1+\phi^2\kappa_2}{\phi^2\kappa_1+\phi^{-2}\kappa_2}=1.307$ for $C_1^{\pm1}$ against $C_1^{\pm2}$ at $D=47$. Numerically (\texttt{linnik\_class\_h5.py}, $N\le4\cdot10^6$, $\kappa_j$ the mean of $c_j^2$; the diagonals are $\kappa_j\pi(N)$ to $3\cdot10^{-4}$):
\begin{center}\footnotesize
\begin{tabular}{r|cc|cc|cc}
\toprule
& \multicolumn{2}{c|}{$D=47$, $\kappa=0.6146,\ 0.8820$} & \multicolumn{2}{c|}{$D=79$, $\kappa=0.5043,\ 0.8965$} & \multicolumn{2}{c}{$\mathrm{corr}(S_1,S_2)$}\\
$N$ & $\frac{E\,S_1^2}{\kappa_1\pi}$ & $\frac{E\,S_2^2}{\kappa_2\pi}$ & $\frac{E\,S_1^2}{\kappa_1\pi}$ & $\frac{E\,S_2^2}{\kappa_2\pi}$ & $47$ & $79$\\
\midrule
$[10^5,2\cdot10^5)$ & $0.421$ & $0.421$ & $0.439$ & $0.434$ & $0.026$ & $0.038$\\
$[3\cdot10^5,6\cdot10^5)$ & $0.425$ & $0.423$ & $0.442$ & $0.438$ & $0.025$ & $0.036$\\
$[10^6,2\cdot10^6)$ & $0.433$ & $0.431$ & $0.448$ & $0.443$ & $0.021$ & $0.029$\\
$[2\cdot10^6,4\cdot10^6)$ & $0.435$ & $0.434$ & $0.450$ & $0.446$ & $0.019$ & $0.029$\\
\bottomrule
\end{tabular}
\end{center}
Both modes show the same deficit as for $h=3$, independently of the size of $\kappa_j$, and slightly less of it for the larger $D$ (the depleted arcs of Lemma~\ref{lem:depletion} are $q^2\ll X/D$). The cross diagonal is below $10^{-4}\pi(N)$; the correlation of $S_1$ and $S_2$ is small and positive and decreases slowly, consistent with $o(1)$ at a logarithmic rate but not decisive. The mean squares of the three class deviations agree with the uncorrelated combination $\frac4{25}(a_1^2E\,S_1^2+a_2^2E\,S_2^2)$ to within $1.5$--$4\%$ (the remainder being the covariance), and the measured ratio for $C_1^{\pm1}$ against $C_1^{\pm2}$ at $D=47$ is $1.309$.

\subsection{Sieve weights: towards Chen's theorem in a single class}\label{subsec:sieve}
For $N=p+Q_C(x,y)$ with $Q_C(x,y)$ an almost prime one sifts $m=N-p$ with weights $\lambda_d$, $d\mid P(z)$, and the class decomposition \eqref{eq:dec} applies term by term:
\begin{gather*}
S^\lambda_C(N):=\sum_{p<N}\Bigl(\sum_{d\mid N-p}\lambda_d\Bigr)r_C(N-p)=S^\lambda_\gen(N)+\frac wh\sum_{\psi\ne1}\bar\psi(C)\,S^\lambda_\psi(N),\\
S^\lambda_\psi(N)=\sum_{p<N}\Bigl(\sum_{d\mid N-p}\lambda_d\Bigr)c_\psi(N-p).
\end{gather*}
The genus part is the sieve problem for the sequence $\{N-p\}$ with the divisor-type weight $\frac wh(1*\chi)$, in which the class group plays no role; the class deviation is again a cusp-form sum, and the Hecke relations make it a combination of dilates of $\vt_\psi$. Throughout, $\lambda_d$ is supported on squarefree $d\le X^\theta$ with $(d,D)=1$ and $|\lambda_d|\le1$ (the Rosser--Iwaniec weights and the weights of the fundamental lemma are of this kind).

\begin{lemma}[Dilation]\label{lem:dilation}
For squarefree $d$ with $(d,D)=1$,
\[
\sum_{k\ge1}c_\psi(dk)\,e(kz)=\sum_{s\mid d}\mu(s)\chi(s)\,c_\psi(d/s)\,\vt_\psi(sz).
\]
\end{lemma}
\begin{proof}
For a prime $\ell\nmid D$ the Hecke relation $c_\psi(\ell k)=c_\psi(\ell)c_\psi(k)-\chi(\ell)c_\psi(k/\ell)$ (with $c_\psi(k/\ell)=0$ if $\ell\nmid k$) reads $U_\ell\vt_\psi=c_\psi(\ell)\vt_\psi-\chi(\ell)\vt_\psi(\ell z)$, and $U_\ell$ commutes with $f(z)\mapsto f(\ell'z)$ for $\ell'\ne\ell$. Applying $U_\ell$ for the primes $\ell\mid d$ in turn and using $c_\psi(d/s)=\prod_{\ell\mid d/s}c_\psi(\ell)$ gives the formula.
\end{proof}

\begin{proposition}[Wilton bound for the dilates]\label{prop:wiltond}
For squarefree $d$ with $(d,D)=1$, uniformly in $\alpha$ and $Y\ge2$,
\begin{gather*}
F_d(\alpha;Y):=\sum_{k\le Y/d}c_\psi(dk)\,e(dk\alpha)\ll_D d^{-1/2+\e}\sqrt Y\log Y,\\
\sum_{k}c_\psi(dk)e^{-dk/Y}e(dk\alpha)\ll_Dd^{-1/2+\e}\sqrt Y .
\end{gather*}
\end{proposition}
\begin{proof}
By Lemma~\ref{lem:dilation}, $F_d(\alpha;Y)=\sum_{s\mid d}\mu(s)\chi(s)c_\psi(d/s)\sum_{j\le Y/(ds)}c_\psi(j)e(jds\alpha)$, each inner sum is $\ll(Y/ds)^{1/2}\log Y$ by Lemma~\ref{lem:wilton} (uniformly in the frequency $ds\alpha$), and $|c_\psi(n)|\le\tau(n)$, $\sum_{s\mid d}1\ll d^\e$. The same for the smooth sums with the second statement of Lemma~\ref{lem:wilton}.
\end{proof}

\begin{theorem}\label{thm:sifted}
Let $\lambda_d$ be as above, of level $X^\theta$, $0\le\theta<1$. Then for every $\psi\ne1$
\[
\sum_{N\le X}|S^\lambda_\psi(N)|^2\ll_{D,\e}X^{2+\theta+\e},
\]
and for every $\eta>0$ all $N\le X$ outside a set of $O(X^{1-2\eta+\e})$ elements satisfy, for every class $C$,
\[
S^\lambda_C(N)=S^\lambda_\gen(N)+O(N^{1/2+\theta/2+\eta}).
\]
\end{theorem}
\begin{proof}
As in Theorem~\ref{thm:almostall}, $S^\lambda_\psi(N)$ is the coefficient of $e(N\alpha)$ in $P(\alpha)F^\lambda(\alpha)$ with
\[
F^\lambda(\alpha)=\sum_{m\le X}\Bigl(\sum_{d\mid m}\lambda_d\Bigr)c_\psi(m)e(m\alpha)=\sum_d\lambda_dF_d(\alpha;X),
\]
so by Parseval and Proposition~\ref{prop:wiltond}
\[
\sum_{N\le X}|S^\lambda_\psi(N)|^2\le\|F^\lambda\|_\infty^2\,\pi(X)\ll\Bigl(\sum_{d\le X^\theta}d^{-1/2+\e}\Bigr)^2X(\log X)^2\,\pi(X)\ll X^{2+\theta+3\e}.
\]
The exceptional set is counted dyadically as in Theorem~\ref{thm:almostall}.
\end{proof}

Since the main term of a sieve of dimension one on $\{N-p\}$ is of size $N/(\log N)^2$, the class deviation is negligible for almost all $N$ for \emph{every} sieve of level $\theta<1$: whatever a sieve proves about $N=p+Q(x,y)$ with $Q(x,y)$ in the genus (an upper bound, a lower bound, the Chen-type combination of sieve functionals) holds for each class separately, for almost all $N$, with the main term divided by $h$. Two things remain. The genus problem itself: for almost all $N$ the input is a level of distribution of the primes with the weight $(1*\chi)(N-p)$ on average over $N$, which we have not worked out (for individual $N$ it is Hypothesis~$\mathrm H(N)$, Proposition~\ref{prop:genus}); for the mere existence of representations with $Q_C(x,y)$ an almost prime this input turns out to be unnecessary for almost all $N$ (Theorem~\ref{thm:goldbach} and Corollary~\ref{cor:odd}). And the switching step of Chen's argument, which counts $N-m$ prime for $m=q_1q_2q_3$: there the class weight $c_\psi(m)$ sits on the almost primes and the sieve on $N-m$, the roles of $P$ and $F$ are exchanged, and the Parseval argument needs the exponential sum of $c_\psi$ over the primes, uniformly in $\alpha$. For holomorphic forms of full level Fouvry and Ganguly \cite{FG} prove $\sum_{n\le X}\Lambda(n)\lambda_f(n)e(n\alpha)\ll X\exp(-c\sqrt{\log X})$. For $\vt_\psi$ the classical method gives the saving $(\log X)^{-A}$, which is all the second-moment argument needs; the one point to check is that the dihedral form, although its twisted $L$-functions are $L$-functions of characters of $K$, has no Siegel zeros, because $h$ is odd.

\begin{proposition}[Vinogradov--Siegel--Walfisz bound for $\vt_\psi$]\label{prop:VSW}
For every $\psi\ne1$ and $A>0$, uniformly in $\alpha\in\mathbb R$,
\[
\sum_{n\le X}\Lambda(n)c_\psi(n)e(n\alpha)\ll_{D,A}X(\log X)^{-A},\qquad \sum_{p\le X}c_\psi(p)e(p\alpha)\ll_{D,A}X(\log X)^{-A}.
\]
\end{proposition}
\begin{proof}
Let $Q=(\log X)^B$ with $B=B(A)$ large, and by Dirichlet's theorem $\alpha=\frac aq+\beta$ with $(a,q)=1$, $q\le X/Q$, $|\beta|\le Q/(qX)$.

\emph{Major arcs, $q\le Q$.} The prime powers $n=p^k$, $k\ge2$, contribute $O(\sqrt X\log X)$ ($|c_\psi(p^k)|\le k+1$), the $n$ with $(n,q)>1$ contribute $O(\omega(q)(\log X)^2)$, and for $(n,q)=1$ we expand $e(na/q)=\frac1{\varphi(q)}\sum_{\xi\bmod q}\tau(\bar\xi,a)\xi(n)$ with $|\tau(\bar\xi,a)|\le q$. Replacing $\xi$ by the primitive character $\xi^*$ modulo $d\mid q$ that induces it changes $\sum_{n\le t}\Lambda(n)c_\psi(n)\xi(n)$ by $O(\omega(q)(\log t)^2)$. Now
\[
\sum_nc_\psi(n)\xi^*(n)n^{-s}=\sum_{\mathfrak a}\psi(\mathfrak a)\xi^*(N\mathfrak a)(N\mathfrak a)^{-s}=L_K(s,\Psi),\qquad\Psi=\psi\cdot(\xi^*\circ N),
\]
a ray class character of $K$ of conductor dividing $(d)$, so $N\mathfrak f\le d^2\le Q^2$, and $\sum_{n\le t}\Lambda(n)c_\psi(n)\xi^*(n)=\sum_{N\mathfrak a\le t}\Lambda_K(\mathfrak a)\Psi(\mathfrak a)+O(\sqrt t\log t)$, the two sums differing at prime powers only. The character $\Psi$ is not real: for a split prime $\mathfrak p\nmid qD$ with $\psi(\mathfrak p)\ne1$ (such exist by Chebotarev, $\psi\ne1$) one has $N\mathfrak p=N\bar{\mathfrak p}$ and $\psi(\bar{\mathfrak p})=\psi(\mathfrak p)^{-1}$, so $\Psi(\mathfrak p)/\Psi(\bar{\mathfrak p})=\psi(\mathfrak p)^2\ne1$ because $\psi$ has odd order, hence $\Psi^2(\mathfrak p)\ne\Psi^2(\bar{\mathfrak p})$ and $\Psi^2\ne1$. So $L_K(s,\Psi^2)$ is entire, the argument of de la Vall\'ee Poussin and Landau gives the zero-free region $\sigma>1-c_K/\log(N\mathfrak f(|t|+2))$ for $L_K(s,\Psi)$ without exceptional zero (the exceptional zero of a Hecke $L$-function occurs only for real characters, \cite[Chapter~5]{IK}), and the explicit formula gives, uniformly for $N\mathfrak f\le(\log t)^{3B}$,
\[
\sum_{N\mathfrak a\le t}\Lambda_K(\mathfrak a)\Psi(\mathfrak a)\ll_{D,B}t\exp(-c\sqrt{\log t}):
\]
the Siegel--Walfisz theorem for $K$, here effective. Partial summation against $e(n\beta)$ costs the factor $1+2\pi|\beta|X\le1+2\pi Q$ (the range $t\le\sqrt X$, where the uniformity in the conductor may fail, contributes $O(|\beta|X)=O(Q)$ trivially), and the sum over $\xi$ costs $q\le Q$: the major arcs give $X\exp(-c'\sqrt{\log X})$.

\emph{Minor arcs, $Q<q\le X/Q$, $|\beta|\le1/q^2$.} Vaughan's identity with $U=V=X^{1/3}$ \cite[Chapter~24]{Davenport} writes the sum as $O(\log X)$ sums of the two types
\begin{gather*}
S_I=\sum_{m\le X^{2/3}}a_m\sum_{n\le X/m}c_\psi(mn)e(mn\alpha)(\log n)^j,\\ S_{II}=\sum_{M<m\le2M}\ \sum_{V<n\le X/m}a_mb_n\,c_\psi(mn)e(mn\alpha),
\end{gather*}
$j\in\{0,1\}$, $X^{1/3}\le M\le X^{2/3}$, $|a_m|\le\tau(m)\log X$, $|b_n|\le\log X$. The Hecke relation $c_\psi(m)c_\psi(n)=\sum_{g\mid(m,n)}\chi(g)c_\psi(mn/g^2)$ inverts to
\begin{equation}\label{eq:hecke}
c_\psi(mn)=\sum_{g\mid(m,n)}\mu(g)\chi(g)\,c_\psi(m/g)\,c_\psi(n/g)\qquad(m,n\ge1).
\end{equation}
In $S_I$ put $m=gm'$, $n=gn'$: the inner sum $\sum_{n'\le X/(g^2m')}c_\psi(n')e(g^2m'n'\alpha)(\log gn')^j$ is $\ll(X/(g^2m'))^{1/2}(\log X)^2$ by Lemma~\ref{lem:wilton} and partial summation, uniformly in the frequency, so with $|c_\psi(n)|\le\tau(n)$
\[
S_I\ll X^{1/2}(\log X)^3\sum_g\frac{\tau(g)}g\sum_{m'\le X^{2/3}/g}\frac{\tau(m')^2}{m'^{1/2}}\ll X^{5/6}(\log X)^C.
\]
In $S_{II}$ put $m=gm'$, $n=gn'$. The terms with $g>G_0=(\log X)^{A+C}$ are bounded trivially: there are $\ll(X/g^2)\log X$ pairs $(m,n)$ with $g\mid(m,n)$, $mn\le X$, and their weights are $\ll\tau(g)\tau(m')\tau(n')(\log X)^2$, so they contribute $\ll X(\log X)^{C}\sum_{g>G_0}\tau(g)g^{-2}\ll X(\log X)^{-A}$. For fixed $g\le G_0$ the sum is $\sum_{m'}\sum_{n'}A_{m'}B_{n'}e(g^2m'n'\alpha)$ with $A_{m'}=a_{gm'}c_\psi(m')$, $M/g<m'\le2M/g$, $B_{n'}=b_{gn'}c_\psi(n')$, $n'\le X/(g^2m')<X/(gM)$, and $\|A\|_2^2\ll\tau(g)^2(M/g)(\log X)^C$, $\|B\|_2^2\ll\tau(g)^2(X/(gM))(\log X)^C$ (from $\sum_{n\le Y}\tau(n)^4\ll Y(\log Y)^{15}$). Cauchy--Schwarz in $n'$ and the geometric progressions in $n'$ give
\begin{multline*}
|S_{II,g}|^2\le\|B\|_2^2\sum_{m_1',m_2'}|A_{m_1'}A_{m_2'}|\min\Bigl(\frac X{gM},\frac1{\|g^2(m_1'-m_2')\alpha\|}\Bigr)\\
\le\|A\|_2^2\|B\|_2^2\sum_{|h|\le gM}\min\Bigl(\frac X{gM},\frac1{\|h\alpha\|}\Bigr),
\end{multline*}
where $h=g^2(m_1'-m_2')$ runs over multiples of $g^2$ with $|h|\le gM$ and we drop the divisibility. The last sum is $\ll(gM/q+1)(X/(gM)+q\log q)\ll(X/q+X/(gM)+gM+q)\log X$ \cite[Chapter~25]{Davenport}, and with $Q<q\le X/Q$, $X^{1/3}\le M\le X^{2/3}$, $g\le G_0$,
\[
|S_{II,g}|\ll\tau(g)^2\frac{X^{1/2}}g(\log X)^C\Bigl(\frac XQ+G_0X^{2/3}\Bigr)^{1/2}\ll\tau(g)^2\frac Xg(\log X)^CQ^{-1/2}+X^{5/6}(\log X)^{C'};
\]
summing over $g\le G_0$ and the $O(\log X)$ values of $M$ gives $S_{II}\ll X(\log X)^{C''}Q^{-1/2}+X^{5/6+\e}$, and $B=2A+2C''$ makes the minor arcs $\ll X(\log X)^{-A}$. The sum over the primes follows by partial summation.
\end{proof}

\begin{corollary}[Strata and the switching set]\label{cor:strata}
Fix $k\ge1$, $\delta>0$ and $A>0$. Uniformly in $\alpha$,
\[
\sum_{m\le X,\ \Omega(m)=k}c_\psi(m)e(m\alpha)\ll_{D,k,A}X(\log X)^{-A},\qquad\sum_{m\in E}c_\psi(m)e(m\alpha)\ll_{D,\delta,A}X(\log X)^{-A}
\]
for every set $E=E(X)$ of products $m=q_1q_2q_3\le X$ of primes $q_1<q_2\le q_3$ with $q_1\in I_1$, $q_2\in I_2$ for intervals $I_1,I_2$ such that $q_1q_2\le X^{1-\delta}$ (Chen's switching set has $X^{1/10}\le q_1<X^{1/3}\le q_2$).
\end{corollary}
\begin{proof}
Write $m=m_0q$ with $q=P^+(m)$ the largest prime factor. If $q\nmid m_0$ then $c_\psi(m)=c_\psi(m_0)c_\psi(q)$ by \eqref{eq:hecke}, and the sum over these $m$ is
\[
\sum_{m_0}c_\psi(m_0)\sum_{P^+(m_0)<q\le X/m_0}c_\psi(q)e(m_0q\alpha),
\]
over $m_0$ with $\Omega(m_0)=k-1$ (resp.\ $m_0=q_1q_2$, $q_1\in I_1$, $q_2\in I_2$). Since $q\ge P^+(m_0)\ge m_0^{1/(k-1)}$ and $m_0q\le X$, $m_0\le X^{(k-1)/k}$, so the inner sum runs over the primes up to $Y=X/m_0\ge X^{1/k}$ (resp.\ $\ge X^\delta$), and Proposition~\ref{prop:VSW} bounds it by $Y(\log Y)^{-A}\ll(X/m_0)(\log X)^{-A}$, while $\sum_{\Omega(m_0)=k-1}|c_\psi(m_0)|/m_0\le2^{k-1}(\sum_{p\le X}1/p)^{k-1}\ll(\log\log X)^{k-1}$. If $q\mid m_0$ then $m=q^2m_1$ with $\Omega(m_1)=k-2$ and $P^+(m_1)\le q$; the number of such $m\le X$ is at most $\sum_q\min(\pi(q)^{k-2},X/q^2)\ll X^{1-1/k}$ (split at $q=X^{1/k}$), and they are bounded trivially, $|c_\psi(m)|\le\tau(m)\le3\cdot2^{k-2}$.
\end{proof}

\begin{theorem}[The strata and the switching sum for almost all $N$]\label{thm:strata}
Fix $k\ge1$ and $A>0$. Let $S^{(k)}_\psi(N)=\sum_{p<N,\ \Omega(N-p)=k}c_\psi(N-p)$, and for weights $w$ with $\sum_{n\le X}w(n)^2\ll X(\log X)^C$ (the indicator of the primes, or $w(n)=\sum_{d\mid n}\lambda_d$ with $|\lambda_d|\le1$) let $S^E_\psi(N)=\sum_{m\in E,\ m<N}c_\psi(m)w(N-m)$ with $E$ as in Corollary~\ref{cor:strata}. Then
\[
\sum_{N\le X}|S^{(k)}_\psi(N)|^2\ll X^3(\log X)^{-A},\qquad\sum_{N\le X}|S^E_\psi(N)|^2\ll X^3(\log X)^{-A},
\]
so that all $N\le X$ outside $O(X(\log X)^{-A})$ exceptions satisfy $|S^{(k)}_\psi(N)|\le N(\log N)^{-A}$ and also $|S^E_\psi(N)|\le N(\log N)^{-A}$. In particular the representations $N=p+Q_C(x,y)$ with $\Omega(Q_C(x,y))=k$ are equidistributed among the classes, up to $N(\log N)^{-A}$, for almost all $N$, and so is the class weight in Chen's switching step.
\end{theorem}
\begin{proof}
Parseval as in Theorem~\ref{thm:almostall} with $F^{(k)}(\alpha)=\sum_{m\le X,\,\Omega(m)=k}c_\psi(m)e(m\alpha)$ and Corollary~\ref{cor:strata}:
\[
\sum_{N\le X}|S^{(k)}_\psi(N)|^2\le\|F^{(k)}\|_\infty^2\int_0^1|P|^2\ll X^2(\log X)^{-2A}\pi(X);
\]
for $S^E_\psi$ the same with $F^E$ and $\int_0^1|\sum_nw(n)e(n\alpha)|^2=\sum_nw(n)^2$. The exceptional sets are counted dyadically.
\end{proof}

Theorem~\ref{thm:strata} explains the shares $\frac13$ of Table~\ref{tab:strata} and closes the class side of Chen's argument: for almost all $N$, Chen's theorem for a single class (every large $N$ is $p+Q_C(x,y)$ with $Q_C(x,y)$ a $P_2$) is equivalent to Chen's theorem for the genus, i.e.\ for the weight $(1*\chi)(N-p)$; in fact Theorem~\ref{thm:goldbach} and Corollary~\ref{cor:odd} below prove the single-class statement directly for almost all $N$, with $Q_C(x,y)$ a prime when $N$ is even. The stratum $k=1$ is the binary Goldbach problem with one prime in a class, and here the genus part needs no hypothesis: on primes $q\ne D$ the genus weight $\frac wh(1*\chi)(q)=\frac2h(1+\chi(q))$ is a character modulo $D$, not a divisor sum.

\begin{theorem}[Goldbach in a single ideal class, almost all $N$]\label{thm:goldbach}
Let $G_C(N)=\sum_{p+q=N}r_C(q)$ (ordered pairs of primes). For every $A>0$, all even $N\le X$ outside $O_{D,A}(X(\log X)^{-A})$ exceptions satisfy, for every class $C$,
\[
G_C(N)=\frac2h\Bigl(1-\frac{\chi(N)}{D-2}\Bigr)\sS_2(N)\,\mathrm{li}_2(N)+O(N(\log N)^{-A}),
\]
with $\sS_2(N)=2\prod_{\ell>2}(1-(\ell-1)^{-2})\prod_{\ell\mid N,\,\ell>2}\frac{\ell-1}{\ell-2}$ and $\mathrm{li}_2(N)=\int_2^{N-2}\frac{\dd t}{\log t\log(N-t)}\sim N(\log N)^{-2}$. In particular almost every even $N$ is $p+q$ with $q$ a prime represented by $Q_C$, for every class $C$, and the $h$ classes share the Goldbach pairs with $q$ split in $K$ equally.
\end{theorem}
\begin{proof}
By \eqref{eq:dec}, $G_C=G_\gen+\frac wh\sum_{\psi\ne1}\bar\psi(C)S^{(1)}_\psi$ with $G_\gen(N)=\frac2h\sum_{p+q=N}(1+\chi(q))+O(1)$ (the term $q=D$), and $S^{(1)}_\psi(N)\le N(\log N)^{-A}$ outside the exceptional set of Theorem~\ref{thm:strata}. For $\sum_{p+q=N}1$ the asymptotic formula $\sS_2(N)\mathrm{li}_2(N)+O(N(\log N)^{-A})$ for all even $N\le X$ but $O(X(\log X)^{-A})$ is the theorem of van der Corput, Chudakov and Estermann, \cite[Theorem~3.7]{Vaughan}. For $\sum_{p+q=N}\chi(q)$ the same proof applies (Siegel--Walfisz for the characters modulo $qD$ on the major arcs, Parseval and Vinogradov's bound for $\sum_p\chi(p)e(p\alpha)$ on the minor arcs); only the major-arc main term changes, through the local factor at $D$: the weight $1+\chi(q)$ summed over the admissible residues $q\not\equiv0,N\pmod D$ gives $\sum_{b\ne0,N}(1+\chi(b))=D-2-\chi(N)$ in place of $D-2$ for $D\nmid N$, and $D-1$ in both cases for $D\mid N$. So $\sum_{p+q=N}(1+\chi(q))=(1-\frac{\chi(N)}{D-2})\sS_2(N)\mathrm{li}_2(N)+O(N(\log N)^{-A})$ outside the exceptional set.
\end{proof}

For odd $N$ the same argument gives the representations with the least number of prime factors that the $2$-adic constraint allows.

\begin{corollary}[Odd $N$]\label{cor:odd}
Let $N$ be odd. For $p>2$ the number $m=N-p$ is even, so $\Omega(m)\ge2$; if $\chi(2)=-1$, i.e.\ $D\equiv3\pmod8$, then $r_C(m)>0$ forces $4\mid m$ and $\Omega(m)\ge3$. Put $J_a(N)=\int_2^{(N-2)/a}\frac{\dd u}{\log u\,\log(N-au)}\sim\frac N{a(\log N)^2}$. For every $A>0$, all odd $N\le X$ outside $O_{D,A}(X(\log X)^{-A})$ exceptions satisfy, for every class $C$,
\begin{align*}
\sum_{p+2q=N}r_C(2q)&=\frac4h\Bigl(1-\frac{\chi(N)}{D-2}\Bigr)\sS_2(N)J_2(N)+O(N(\log N)^{-A})\qquad(D\equiv7\bmod8),\\
\sum_{p+4q=N}r_C(4q)&=\frac2h\Bigl(1-\frac{\chi(N)}{D-2}\Bigr)\sS_2(N)J_4(N)+O(N(\log N)^{-A})\qquad(D\equiv3\bmod8).
\end{align*}
In particular almost every odd $N$ is $p+Q_C(x,y)$ with $Q_C(x,y)=2q$, resp.\ $4q$, $q$ prime, for every class $C$. Together with Theorem~\ref{thm:goldbach}: for almost all $N$ and every class $C$, $N=p+Q_C(x,y)$ with $Q_C(x,y)$ a prime ($N$ even), a $P_2$ ($N$ odd, $D\equiv7\bmod8$) or a $P_3$ ($N$ odd, $D\equiv3\bmod8$), unconditionally.
\end{corollary}
\begin{proof}
Let $D\equiv7\pmod8$. For an odd prime $q\ne D$ the ideals of norm $2q$ are $\mathfrak p_2^{\pm}\mathfrak q^{\pm}$, $(1*\chi)(2q)=2(1+\chi(q))$ and $c_\psi(2q)=c_\psi(2)c_\psi(q)$, so by \eqref{eq:dec}
\[
\sum_{p+2q=N}r_C(2q)=\frac4h\sum_{p+2q=N}(1+\chi(q))+\frac wh\sum_{\psi\ne1}\bar\psi(C)\,c_\psi(2)\sum_{p+2q=N}c_\psi(q)+O(1).
\]
The last sums are the coefficients of $P(\alpha)F(2\alpha)$ with $F(\alpha)=\sum_{q\le X/2}c_\psi(q)e(q\alpha)$, and Proposition~\ref{prop:VSW} at the frequency $2\alpha$ with Parseval gives $\sum_{N\le X}|\sum_{p+2q=N}c_\psi(q)|^2\ll X^3(\log X)^{-A}$ as in Theorem~\ref{thm:strata}. For $\sum_{p+2q=N}(1+\chi(q))$ the proof of Theorem~\ref{thm:goldbach} applies to the equation $p+2q=N$: the local factors are those of $\sS_2(N)$ (at $\ell=2$ the one admissible pair $(1,1)$ satisfies $1+2\equiv N$; at odd $\ell$ the pairs $b_1+2b_2\equiv N$ with $b_1b_2\not\equiv0$ number $\ell-2$ or $\ell-1$ according as $\ell\nmid N$ or $\ell\mid N$, as for $b_1+b_2\equiv N$), at $D$ the weight $1+\chi(q)$ summed over $q\not\equiv0,N/2$ gives $D-2-\chi(N/2)=D-2-\chi(N)$ because $\chi(2)=1$, and the integral is $J_2(N)$. For $D\equiv3\pmod8$ the ideal $(2)$ is a principal prime of norm $4$, so $r_C(4q)=r_C(q)$, and the same argument with $p+4q=N$, $\chi(4)=1$ and $J_4(N)$ applies.
\end{proof}

Numerically ($D=23$ and $31$, both $\equiv7\bmod8$; $4000$ random odd $N$ per range): $\#\{p+2q=N\}$ is $0.9919$, $0.9953$, $0.9975$, $0.9983$ of $\sS_2(N)J_2(N)$ in the four ranges of Table~\ref{tab:strata}, and $\sum_{p+2q=N}r_C(2q)$ against the main term of Corollary~\ref{cor:odd} is $0.992$ ($C_0$), $0.984$ ($C_1$) at $N\approx1.5\cdot10^5$ and $0.998$, $0.997$ at $N\approx3\cdot10^6$ for $D=23$, the same to $\pm0.002$ for $\chi(N)=\pm1$ and for $D=31$. Here it is $C_1$ that lags: the Chebyshev bias of the principal primes enters through $c_\psi(2)c_\psi(q)$ with $c_\psi(2)=-1$, which turns the deficit of $G_{C_0}$ in Theorem~\ref{thm:goldbach} into a deficit of the non-principal class.

Numerically (\texttt{goldbach\_class\_parity.py}, $4000$ random even $N$ in each range, $\sS_2(N)\mathrm{li}_2(N)$ by quadrature): for $D=23$ the full Goldbach count $G(N)$ is $0.9933$, $0.9961$, $0.9979$, $0.9985$ of $\sS_2(N)\mathrm{li}_2(N)$ in the four ranges of Table~\ref{tab:strata} (standard error $0.0002$), and $G_{C_0}(N)$, $G_{C_1}(N)$ against the main term of Theorem~\ref{thm:goldbach} are $0.985$, $0.992$ at $N\approx1.5\cdot10^5$ and $0.996$, $0.999$ at $N\approx3\cdot10^6$ for $\chi(N)=1$, the same to $\pm0.001$ for $\chi(N)=-1$; the deficit of $G_{C_0}$ against $G_{C_1}$ is the Chebyshev bias of the principal primes of Remark~\ref{rem:parity}, and it decreases. The local factor at $D$ is tested by the ratio of the split-pair counts $\frac14(G_{C_0}+2G_{C_1})$ for $\chi(N)=1$ against $\chi(N)=-1$, which should be $\frac{D-3}{D-1}$: measured $0.9081$, $0.9096$, $0.9089$, $0.9093$ against $0.9091$ for $D=23$ (the alternative factor $1-\chi(N)/(D-1)$, which is the one of the Linnik problem, Proposition~\ref{prop:genus}, would give $0.9130$), and $0.9340$, $0.9343$, $0.9334$, $0.9332$ against $0.9333$ for $D=31$ ($0.9355$ for the alternative); for $D\mid N$ the split share is $0.500\pm0.001$.

\begin{remark}[No class-group grading of the parity obstruction]\label{rem:parity}
On the integers $m$ all of whose prime factors are split and non-principal one has $c_\psi(m)=(-1)^{\Omega(m)}$ for $h=3$ ($c_\psi(q)=-1$ for such $q$), so that the primes $q$ in this set are counted by $r_{C_1}$ only ($r_{C_0}(q)=0$, $r_{C_1}(q)=2$) and the semiprimes $q_1q_2$ twice as much by $r_{C_0}$ as by $r_{C_1}$ ($4$ against $2$): the ideals $\mathfrak q_1\bar{\mathfrak q}_2$, $\bar{\mathfrak q}_1\mathfrak q_2$ are principal. One might hope that this grading separates the primes from the semiprimes in a single class and so circumvents the parity obstruction of Chen's theorem. It does not: the condition ``non-principal prime factors'' is itself a class-group condition, invisible to a sieve over $\Z$, and on the full stratum $\Omega(m)=k$ the mean of $c_\psi$ is zero for every $k$ (Chebotarev and multiplicativity: $c_\psi(q)=2,-1,0$ with densities $\frac16,\frac13,\frac12$ for $h=3$, and $c_\psi(q_1q_2)=4,-2,1$ with densities $\frac19,\frac49,\frac49$ of the split pairs). The strata $\Omega(N-p)=k$ of $S_{C}(N)$ are therefore equidistributed among the classes like $S_C$ itself (Theorem~\ref{thm:strata}; Table~\ref{tab:strata}: shares within $0.002$ of $\frac13$), and the stratum $\Omega(N-p)=1$, the Goldbach pairs $N=p+q$ with $q=Q_C(x,y)$, is one third of the Goldbach pairs with $q$ split for every class. What the strata do show is a Chebyshev bias of lower order: $\sum_{q\le x}c_\psi(q)\approx-\sqrt x/\log x$ by the explicit formula (the prime squares, $c_\psi(q^2)=c_\psi(q)^2-\chi(q)$ with mean $1$, enter $\sum_{n\le x}\Lambda(n)c_\psi(n)$ with a positive sign), the squares $q^2$ and the products with the small non-principal primes ($c_\psi(2)=c_\psi(3)=-1$ for $D=23$) propagate it with alternating signs to the strata, and the bias $E\,S^{(k)}_\psi/\mathrm{rms}\,S^{(k)}_\psi$ is negative, positive, negative for $k=1,2,3$ (Table~\ref{tab:strata}: $-0.1$, $+0.3$, $-0.5$ at $N\approx10^5$), of relative order $(\log N)^{-3/2}$; it decreases slowly for $k=2,3$ and is flat within the noise for $k=1$ at these sizes. In the shares it is visible as $0.332$--$0.333$ for $C_0$ on the Goldbach stratum (fewer principal primes) and $0.334$--$0.335$ on the semiprimes.
\end{remark}

\begin{proposition}[Major arcs under sieve weights]\label{prop:siftedarcs}
Let $\lambda_d$ be of level $X^\theta$ and let $\mathfrak M_\theta$ be the arcs of Section~\ref{subsec:reduction} with $Q_\theta=X^{1/2-\theta}(\log X)^{-B}$ in place of $Q$; put
\[
F^\lambda(\alpha)=\sum_d\lambda_d\sum_kc_\psi(dk)e^{-dk/X}e(dk\alpha).
\] Then $\sup_{\mathfrak M_\theta}|F^\lambda|\ll_AX^{-A}$.
\end{proposition}
\begin{proof}
By Lemma~\ref{lem:dilation}, $F^\lambda$ is a sum of $\ll X^{\theta+\e}$ terms $\vt_\psi(sd(\alpha+\frac i{2\pi X}))$ with $sd\le d^2\le X^{2\theta}$ and bounded coefficients. With $X'=X/(sd)$, $\alpha'=sd\alpha=\frac{a'}{q'}+\beta'$, $q'\le q$, $|\beta'|\le sdK/X=K/X'$, the proof of Theorem~\ref{thm:reduction} applies to $\vt_\psi(\alpha'+\frac i{2\pi X'})$ as soon as $X'/q'^2\ge X/(X^{2\theta}q^2)\ge(\log X)^{2B}$, which is $q\le Q_\theta$.
\end{proof}

So a sieve of level $X^\theta$ keeps the depletion of Lemma~\ref{lem:depletion} on the arcs of order $X^{1/2-\theta}$, and the reduction of Conjecture~\ref{conj:var} carries over to $S^\lambda_\psi$ with the arcs $\mathfrak M_\theta$, which carry the mass $(\frac12-\theta)\int|P|^2$ by Hardy--Littlewood. Beyond $Q_\theta$ the dilates $\vt_\psi(nz)$ with $n\ge X/q^2$ are no longer small at $a/q$, and their share of $\|F^\lambda\|_2^2$ (the dilate $\vt_\psi(nz)$ carries $\asymp\|\vt_\psi\|_2^2/n$) is what the major arcs see. For the sharp sieve $(m,P(z))=1$, $z=X^{1/\beta}$, i.e.\ $\lambda_d=\mu(d)$ on $d\mid P(z)$ without truncation, the weight $\sum_{d\mid P(z),\,d\ge T}1/d$ of the undepleted dilates at $a/q$, $T=X/q^2$, is a positive proportion of $\sum_{d\mid P(z)}1/d=\prod_{\ell<z}(1+\ell^{-1})$ whenever $\log T\ll\log z$, so the variance ratio should exceed $\frac12$ by an amount that grows with the sifting, $\frac12+O(1/\beta)$ in the limit, and tend to $1$ for the strata $\Omega(m)=k$, which have no modular structure at all (there $F$ at $a/q$ is governed by the distribution of $c_\psi$ over the primes in residue classes modulo $q$, a square-root cancellation under GRH, the same order as on the minor arcs). The data confirm the ordering (Table~\ref{tab:strata}, script \texttt{goldbach\_class\_parity.py}): in every range and for both $D$ the ratio rises monotonically from the unsifted $0.41$--$0.43$ through $\beta=6,4,3$ to the strata $\Omega=1,2,3$. The sifted sets with $\beta=3$ ($N-p$ free of prime factors below $N^{1/3}$: the Chen set, $\Omega(N-p)\le2$ up to a few $\Omega=3$) are at $0.45$--$0.50$, i.e.\ their class deviation has lost about half of the depletion at these sizes, and the strata are at $0.54$--$0.81$, well above the unsifted value but still far from $1$; the ratios for $\Omega=1,2$ increase with $N$, the one for $\Omega=3$ does not yet. The classes stay equidistributed in every sifted set and every stratum (shares within $0.002$ of $\frac13$, the deviation being the Chebyshev bias of Remark~\ref{rem:parity}), and the sieve adds no bias to the class deviation (the bias of the sifted sets is $\le0.04$ of the rms, the strata carry the alternating Chebyshev bias).

\begin{table}[htbp]
\centering\footnotesize\setlength{\tabcolsep}{3.2pt}
\begin{tabular}{r|ccccccc|ccccccc}
\toprule
& \multicolumn{7}{c|}{$D=23$: $E\,(S^u_c)^2/E\,\mathrm{Diag}^u$} & \multicolumn{7}{c}{$D=31$}\\
$N$ from & $u=1$ & $\beta=6$ & $\beta=4$ & $\beta=3$ & $\Omega=1$ & $\Omega=2$ & $\Omega=3$ & $u=1$ & $\beta=6$ & $\beta=4$ & $\beta=3$ & $\Omega=1$ & $\Omega=2$ & $\Omega=3$\\
\midrule
$10^5$ & $0.405$ & $0.414$ & $0.428$ & $0.446$ & $0.541$ & $0.616$ & $0.813$ & $0.414$ & $0.430$ & $0.448$ & $0.466$ & $0.576$ & $0.637$ & $0.780$\\
$3\cdot10^5$ & $0.414$ & $0.422$ & $0.436$ & $0.466$ & $0.554$ & $0.636$ & $0.814$ & $0.422$ & $0.439$ & $0.449$ & $0.477$ & $0.575$ & $0.645$ & $0.813$\\
$10^6$ & $0.423$ & $0.431$ & $0.454$ & $0.476$ & $0.587$ & $0.647$ & $0.775$ & $0.429$ & $0.447$ & $0.461$ & $0.497$ & $0.596$ & $0.668$ & $0.783$\\
$2\cdot10^6$ & $0.427$ & $0.436$ & $0.457$ & $0.485$ & $0.612$ & $0.679$ & $0.774$ & $0.432$ & $0.449$ & $0.461$ & $0.500$ & $0.590$ & $0.629$ & $0.730$\\
\midrule
\multicolumn{15}{l}{bias $E\,S^u_c/\mathrm{rms}\,S^u_c$ on the strata $\Omega=1,2,3$ ($D=23$; $D=31$):}\\
$10^5$ & \multicolumn{7}{l|}{$-0.119$, $+0.298$, $-0.505$} & \multicolumn{7}{l}{$-0.143$, $+0.296$, $-0.466$}\\
$2\cdot10^6$ & \multicolumn{7}{l|}{$-0.128$, $+0.269$, $-0.411$} & \multicolumn{7}{l}{$-0.087$, $+0.211$, $-0.375$}\\
\bottomrule
\end{tabular}
\medskip
\caption{The variance ratio of the class deviation under sieve weights $u(N-p)$, means over all $N$ in $[N_0,2N_0)$ (\texttt{goldbach\_class\_parity.py}): $u=1$ (Table~\ref{tab:ratios}), $u=1_{(m,P(z))=1}$ with $z=N_0^{1/\beta}$, and the strata $u=1_{\Omega(m)=k}$. The diagonal $E\,\mathrm{Diag}^u=E\sum_pu(N-p)c(N-p)^2$ is $\kappa^u\pi(N)$ with $\kappa^u=0.42$, $0.17$, $0.14$, $0.11$--$0.13$, $0.07$--$0.09$, $0.14$--$0.15$, $0.09$--$0.10$ in the order of the columns ($D=23$). The class shares are within $0.002$ of $\frac13$ in every row. The sifted strata $\beta=3$, $\Omega=1$ and $\beta=3$, $\Omega=2$ (the Goldbach pairs and the Chen semiprimes with large prime factors) are at $0.54$--$0.61$ and $0.64$--$0.69$ for $D=23$.}
\label{tab:strata}
\end{table}

\begin{remark}[Quadratic class weights do not help either]\label{rem:hbt}
Remark~\ref{rem:parity} concerns the linear combinations of the $r_C$. One may ask whether a weight quadratic in the class representation numbers separates the primes from the semiprimes. For $h=3$ the cross-class coincidence $\mathcal C(m)=r_{C_0}(m)r_{C_1}(m)$, the number of simultaneous representations $Q_{C_0}(x,y)=Q_{C_1}(u,v)=m$, vanishes on every prime (the prime ideals above $q$ lie in one form class: $(r_{C_0},r_{C_1})(q)=(4,0)$, $(0,2)$ or $(0,0)$), and on the split semiprimes $q_1q_2$ it is $8$ when both prime ideals are non-principal ($(r_{C_0},r_{C_1})=(4,2)$, density $\frac49$ of the split pairs) and $0$ otherwise ($(8,0)$ and $(0,4)$). So $w_0=r_{C_0}-\frac12\mathcal C$ and $w_1=r_{C_1}-\frac14\mathcal C$ agree with $r_{C_0}$, $r_{C_1}$ on the primes and annihilate the semiprimes with two non-principal factors, which carry $\frac23$ of the semiprime weight of $r_{C_0}$ and $\frac13$ of that of $r_{C_1}$ (for odd $h$ in general, $r_{C_0}-\frac14r_{C_0}(2(1*\chi)-r_{C_0})$ removes the fraction $\frac{h-1}h$). On the Chen set $(N-p,P(N^{1/3}))=1$ this raises the share of the Goldbach primes in the weighted count (\texttt{goldbach\_class\_parity.py}, $D=23$, $N\in[2\cdot10^6,4\cdot10^6)$: the coincidence sum $\sum_p\mathcal C(N-p)$ over the sifted $N-p$ is $0$ on $\Omega=1$ exactly, and on $\Omega=2$ it is $1.42$ times the $r_{C_0}$-sum and, divided by $4$, $0.356$ times the $r_{C_1}$-sum, against $\frac43$ and $\frac13$; the share of the primes in the weighted count rises from $0.61$ to $0.85$ for $C_0$ and from $0.61$ to $0.71$ for $C_1$, and the $\Omega=3$ part of $w_0$ is negative, $-16$ against $+14.5$ for $r_{C_0}$). But this is bookkeeping, not arithmetic. The weight $\mathcal C$ is the coefficient of the weight-two form $\vt_{C_0}\vt_{C_1}=r_\gen^2+\frac23r_\gen c_\psi-\frac89c_\psi^2$ (its three Rankin--Selberg channels are those of Part~VI \cite{ChenVI}), $w_0$ is not non-negative ($w_0=-8$ on the products of three non-principal split primes, which survive the sieve at $N^{1/3}$), and the semiprimes with two principal factors keep the weight $8$. Above all, a sieve sees only the distribution of the weighted sequence in arithmetic progressions, and the parity obstruction of Selberg is a statement about that information: the sequence $\{N-p\}$ with the weight $w_0$ and the same sequence with the weight $w_0(N-p)(1+\lambda(N-p))$, which has no elements with $\Omega$ odd, have the same distribution in progressions to the accuracy a sieve can use, for any weight built from the coefficients of fixed modular forms. The pointwise identities above change which elements are weighted, not what the sieve can know about them.
\end{remark}

\begin{remark}[The dual Titchmarsh input for weight one]\label{rem:ABL}
We checked how far the proof of \cite[Theorem~1.3]{ABL} uses the trivial nebentypus. The form $\varphi$ enters only through (i) the Rankin--Selberg bound $\sum_{n\le X}|\lambda(n)|^2\ll(NX)^\e X$, (ii) the Bessel transforms of their Lemma~7.5, and (iii) the Voronoi formula with additive twists, their Lemma~7.4, quoted from the reference [BMN] of their bibliography ``in the case of trivial nebentypus''; the Kloosterman sums produced by (iii) are then treated by their Theorem~2.4 through the Kuznetsov formula. For $\vt_\psi$, (i) holds, (ii) holds with $J_{k-1}=J_0$ (and the Ramanujan bound holds by Deligne--Serre). For (iii) with the nebentypus $\omega=\chi_{-D}$ and $(q,D)=1$ the computation in the proof of Lemma~\ref{lem:depletion} gives the Voronoi formula with the dual twist $e(-n\overline{aD}/q)$ and the additional factor $\omega(q)$, independent of $a$ but not of $q$, so that the sums of Kloosterman sums in their Section~8 acquire the weight $\omega(c)$: they become Kloosterman sums for the pair of cusps $(\infty,0)$ of $\Gamma_0(D)$ with the character $\omega$, and their Theorem~2.4 has to be re-proved with the Kuznetsov formula for this pair of cusps and nebentypus; for $D\mid q$ the Voronoi coefficients of their Lemma~7.4 depend on $a_1$ through a character and the Kloosterman sums become twisted. The tools for both exist (the Kuznetsov formula for general cusps and characters, the bound $\theta\le7/64$ for every nebentypus), but the extension is not in the literature as far as we know and we have not carried it out; Hypothesis~\ref{hyp:T} stays a hypothesis.
\end{remark}

\section*{Data and code}
{\raggedright \texttt{code/linnik\_class.py} (Table~\ref{tab:ratios}, outputs \texttt{results/linnik\_class\_23.txt}, \texttt{\_31.txt}), \texttt{code/linnik\_class\_variance.py} (mean square, diagonal, the Wilton supremum: $\sup_\alpha|\sum_{m\le Y}c(m)e(m\alpha)|=1.51\sqrt Y$, $1.56\sqrt Y$, $1.58\sqrt Y$ at $Y=10^5,10^6,4\cdot10^6$ for $D=23$), \texttt{code/linnik\_class\_arcs.py} (the depletion profile $\rho(q)$ and the model), \texttt{code/linnik\_class\_smooth.py} (the smoothed variance against the $\beta$-integrated model), \texttt{code/linnik\_class\_split.py} (the major/minor split of Section~\ref{subsec:reduction}), \texttt{code/linnik\_class\_h5.py} (class number five, outputs \texttt{results/linnik\_class\_h5\_47.txt}, \texttt{\_79.txt}), \texttt{code/goldbach\_class\_parity.py} (Section~\ref{subsec:sieve}: sieve weights and $\Omega$-strata, outputs \texttt{results/goldbach\_class\_parity\_23.txt}, \texttt{\_31.txt}). The repository is \url{https://github.com/Ruqing1963/two-squares-linnik-class}; this paper is archived as doi:10.5281/zenodo.24064281.\par}

\begin{thebibliography}{9}\raggedright
\bibitem{ABL} E. Assing, V. Blomer, J. Li, Uniform Titchmarsh divisor problems, arXiv:2005.13915 (2020).
\bibitem{Davenport} H. Davenport, Multiplicative Number Theory, 3rd ed., Graduate Texts in Math. 74, Springer, New York, 2000.
\bibitem{IK} H. Iwaniec, E. Kowalski, Analytic Number Theory, Amer. Math. Soc. Colloq. Publ. 53, Providence, 2004.
\bibitem{Vaughan} R. C. Vaughan, The Hardy--Littlewood Method, 2nd ed., Cambridge Tracts in Math. 125, Cambridge Univ. Press, 1997.
\bibitem{ChenL} R. Chen, The Linnik problem for the norms of the imaginary quadratic fields of class number one, unpublished draft (2026), \url{https://github.com/Ruqing1963/two-squares-linnik-class}, \texttt{paper/linnik\_heegner.pdf}.
\bibitem{ChenV} R. Chen, Sums of two squares and Heegner norms with few prime factors, V: ideal class groups, weight-one dihedral cusp forms, and the splitting of the scalloped norm, preprint (2026), doi:10.5281/zenodo.23923518.
\bibitem{ChenVI} R. Chen, Sums of two squares and Heegner norms with few prime factors, VI: quaternary theta series of ideal classes, dihedral triple products, and a rank-one Gross formula, preprint (2026), doi:10.5281/zenodo.23924216.
\bibitem{ChenVIII} R. Chen, Sums of two squares and Heegner norms with few prime factors, VIII: the mixed form $x^2+y^2+u^2+uv+v^2$, a quaternary form without cusp forms, and the sieve, preprint (2026), doi:10.5281/zenodo.23923846.
\bibitem{Montgomery} H. L. Montgomery, Primes in arithmetic progressions, Michigan Math. J. 17 (1970), 33--39.
\bibitem{Hooley} C. Hooley, On the Barban--Davenport--Halberstam theorem I, J. reine angew. Math. 274/275 (1975), 206--223.
\bibitem{FG} \'E. Fouvry, S. Ganguly, Strong orthogonality between the M\"obius function, additive characters and Fourier coefficients of cusp forms, Compos. Math. 150 (2014), 763--797.
\bibitem{GM} D. A. Goldston, H. L. Montgomery, Pair correlation of zeros and primes in short intervals, in: Analytic Number Theory and Diophantine Problems (Stillwater 1984), Progr. Math. 70, Birkh\"auser, Boston, 1987, 183--203.
\bibitem{Pitt} N. J. E. Pitt, On an analogue of Titchmarsh's divisor problem for holomorphic cusp forms, J. Amer. Math. Soc. 26 (2013), 735--776.
\end{thebibliography}
\end{document}
